\section{Controlled component ablations and uncertainty stress tests}
\label{sec:p18-ablation}
We performed single-component ablations on the same \(90\) seeded synthetic cases used in Section~\ref{sec:p17-benchmark} (seed \(1717\); \(30\) cases each with full, hold and linear timer reset). The reference outcome is the joint, independent-rectangular model criterion under the fixed-template monotonicity assumptions of Section~\ref{sec:p13-continuum}. All screens retain finite relay-assertion requirements. In each ablation only the named constraint or approximation changes; the scenarios, relay parameters, clearing deadlines and frequency data are held fixed.

The complete criterion approved \(22/90\) cases. Omitting the frequency gate yielded \(36\) approvals and \(14\) false approvals. Omitting the primary clearing deadline yielded \(23\) approvals and \(1\) false approval; omitting the counterfactual backup deadline yielded \(28\) approvals and \(6\) false approvals. Omitting the coordination margin yielded \(44\) approvals and \(22\) false approvals. Replacing independent current scales by common-derating endpoint checks yielded \(25\) approvals and \(3\) false approvals. Suppressing uncertain breaker delays by setting both to zero yielded \(22\) approvals and no observed false approvals. These counts concern the chosen synthetic dataset and do not imply that breaker delay can be neglected generally.

For a deliberately simplified timer-memory ablation, progress was discarded at \emph{every} segment boundary, even where the current remained above pickup. This approximation yielded \(20\) approvals, zero false approvals and two conservative rejections. It is not a manufacturer relay reset rule and should not be interpreted as a physically validated no-memory implementation. The experimental distinction is important: a simplified timer may lose useful approvals without necessarily causing unsafe ones in the observed sample.

We also widened the independent current-scaling interval to \([1-w,1]\) while retaining all other case parameters. At \(w=0,0.1,0.2,0.3,0.4\), the respective approval counts were \(40,34,27,24,11\) out of \(90\). This monotone decrease is expected when the admissible uncertainty set expands; it is not an empirical prediction of field restoration rates.

\paragraph{Independent validation and interpretation.}
The CI workflow reran the complete model reference, replayed every reported false-approval witness by deterministic scenario identifier, verified confusion matrices and pairwise overlap counts, checked uncertainty-width monotonicity, and executed repeated wall-clock timing measurements. Seventeen regression and mathematical tests passed. The witness records are the \emph{first observed failing cases} in deterministic identifier order, not formally minimal parameter counterexamples. Timing variability is descriptive only; no statistical significance is asserted. The separate correlated-derating interior counterexample remains an analytic warning that a finite diagonal grid cannot certify a correlated continuum. None of these results establishes installed relay behavior, physical switching feasibility, or a field restoration authorization.